\section*{Chapter \ref{ch:b1:active-margins} --- Active Margins}

\begin{solution}{exo:b1:active-margins:1}
(a) $2\pi G\,\Delta\rho\,h = 2\pi\times\num{6.674e-11}\times1700\times3000 =
\qty{2.14e-3}{m.s^{-2}}$: gravity would be lower by about
\qty{214}{\milli\gal}. (b) For \qty{3.5}{km},
$2\pi\times\num{6.674e-11}\times1700\times3500 = \qty{2.50e-3}{m.s^{-2}}$,
\qty{250}{\milli\gal}: the size of the observed low, so the trench is held far
from isostatic equilibrium. (c) The slab formula assumes the missing rock
extends without limit on every side; a trench only tens of kilometres wide
removes much less mass around the point of measurement, so its water alone
gives less than \qty{250}{\milli\gal}. The rest of the deficit lies deeper:
the whole plate is bent down beneath the trench, and its \omterm{def:b1:layered-earth:layers}{crust}, lighter than
the \omterm{def:b1:layered-earth:layers}{mantle}, takes the place of \omterm{def:b1:layered-earth:layers}{mantle} at its base.
\end{solution}

\begin{solution}{exo:b1:active-margins:2}
(a) $\qty{82}{mm.yr^{-1}} = \qty{82}{km.Myr^{-1}}$, so \qty{820}{km} in
\qty{10}{Myr}. (b) $820\sin\ang{45} = \qty{580}{km}$. (c) The path to
\qty{600}{km} is $600/\sin\ang{45} = \qty{849}{km}$, covered in
$849/82 = \qty{10.3}{Myr}$.
\end{solution}

\begin{solution}{exo:b1:active-margins:3}
$g = 110/\tan\delta$: \qty{191}{km} for \ang{30}, \qty{110}{km} for \ang{45},
\qty{64}{km} for \ang{60}. The arc above the gently dipping slab lies farthest
from its trench.
\end{solution}

\begin{solution}{exo:b1:active-margins:4}
(a) Convergence: $\qty{2.7e6}{m^2.yr^{-1}}/\qty{5.1e7}{m} =
\qty{0.053}{m.yr^{-1}}$, about \qty{53}{mm.yr^{-1}}. Spreading:
$\qty{3.0e6}{m^2.yr^{-1}}/\qty{6.7e7}{m} = \qty{0.045}{m.yr^{-1}}$, about
\qty{45}{mm.yr^{-1}}. Along a shorter total length, the subduction boundaries
consume ocean floor faster than the ridges make it, yet consume less in all.
(b) The other \qty{0.3}{km^2} or so is taken up along the other converging
boundaries of the model, such as continental collisions and convergent zones
within the oceans, where the plates are shortened rather than subducted; and
the ridges are not the only diverging boundaries: \omterm{def:b1:rifts-ridges-margins:rift}{continental rifts} and some
transforms also open a little.
\end{solution}

\begin{solution}{exo:b1:active-margins:5}
(a) $\bar x = 350$, $\bar z = 290$; the deviations of $x$ are $-150$, $-50$,
$50$, $150$ and those of $z$ are $-140$, $-40$, $40$, $140$, so
$S_{xx} = 2(150^2 + 50^2) = \num{50000}$ and
$S_{xz} = 2(150\times140 + 50\times40) = \num{46000}$. Then $s = 0.92$,
$a = 290 - 0.92\times350 = \qty{-32}{km}$, and the \omterm{def:b1:geological-maps:attitude}{dip} is
$\arctan 0.92 = \ang{42.6}$. (b) $z = 0$ at $x = 32/0.92 = \qty{34.8}{km}$
from the trench. (c) Through the end points the slope is $280/300 = 0.933$, a
\omterm{def:b1:geological-maps:attitude}{dip} of \ang{43.0}: close, but it uses only two of the four earthquakes.
\end{solution}

\begin{solution}{exo:b1:active-margins:6}
(a) $50/\tan\ang{25} + 60/\tan\ang{45} = 107.2 + 60.0 = \qty{167.2}{km}$.
(b) A slab of \ang{45} all the way gives \qty{110}{km}: the gentle shallow
segment widens the gap by \qty{57}{km}. (c) The deep segment is
$z = 50 + (x - 107.2)$; it reaches $z = 0$ at $x = \qty{57.2}{km}$, landward of
the trench, as the Kuril line does at \qty{53}{km}: a gentle shallow segment
explains the Kuril offset. (d) $107.2 + 60/\tan\ang{60} = 107.2 + 34.6 =
\qty{141.9}{km}$.
\end{solution}

\begin{solution}{exo:b1:active-margins:7}
(a) The radius is half the \omterm{def:b1:geological-maps:attitude}{dip}, and $r = R\,\theta$ with $\theta$ in radians:
\ang{22.5} and $6371\times0.3927 = \qty{2502}{km}$ for \ang{45}; \ang{15} and
\qty{1668}{km} for \ang{30}; \ang{35} and \qty{3892}{km} for \ang{70}.
(b) The arcs above gently dipping slabs. (c) The prediction is
$44.9/2 = \ang{22.45}$; the fit, \ang{21.9}, is
$0.55/22.45 = \qty{2.4}{\percent}$ smaller.
\end{solution}

\begin{solution}{exo:b1:active-margins:8}
(a) $q/k$: $0.035/2.5 = \qty{0.014}{K.m^{-1}} = \qty{14}{K.km^{-1}}$ beneath
the \omterm{def:b1:active-margins:active-margin}{forearc}, $0.080/2.5 = \qty{32}{K.km^{-1}}$ beneath the back-arc.
(b) $10 + 14\times30 = \qty{430}{\celsius}$ and
$10 + 32\times30 = \qty{970}{\celsius}$. (c) Part of the surface \omterm{def:b1:heat-budget:heat-flow}{heat flow} is
produced by radioactivity within the \omterm{def:b1:layered-earth:layers}{crust} itself, so less heat crosses the
deeper levels and the gradient decreases with depth: a linear extrapolation of
the surface gradient overestimates the temperature at depth.
\end{solution}

\begin{solution}{exo:b1:active-margins:9}
(a) \qty{32}{mm.yr^{-1}} is \qty{32}{km.Myr^{-1}}; the half-rate is
\qty{16}{mm.yr^{-1}}. (b) $100/32 = \qty{3.1}{Myr}$. (c) The rate need not
have been constant; part of the width
may come from the stretching and rifting of the arc before spreading began;
and the axis may have jumped or spread asymmetrically.
\end{solution}

\begin{solution}{exo:b1:active-margins:10}
(a) The path is $600/\sin\ang{45} = \qty{849}{km}$, covered at
\qty{82}{km.Myr^{-1}} in \qty{10.3}{Myr}, that is \qty{3.27e14}{s}.
(b) $\sqrt{\num{e-6}\times\num{3.27e14}} = \qty{1.8e4}{m}$, about
\qty{18}{km}. (c) $(\qty{5e4}{m})^2/\qty{e-6}{m^2.s^{-1}} = \qty{2.5e15}{s}$,
about \qty{79}{Myr}. (d) Heat has reached only the outer \qty{18}{km} of each
side of the slab: its core arrives at \qty{600}{km} almost as cold as it left
the trench, colder and stronger than the \omterm{def:b1:layered-earth:layers}{mantle} around it, which is why it can
still break in earthquakes at that depth.
\end{solution}

\begin{solution}{exo:b1:active-margins:11}
(a) At the distance $L$ behind the toe, the surface stands $L\tan\alpha$ above
the level of the toe and the base lies $L\tan\beta$ below it: the back of the
prism is $L(\tan\alpha + \tan\beta)$ high, and the triangle has the area
$\frac12 L\times L(\tan\alpha + \tan\beta)$. (b) $\tan\ang{4} + \tan\ang{6} =
0.175$. The prism receives $\qty{1}{km}\times\qty{40}{km.Myr^{-1}} =
\qty{40}{km^2.Myr^{-1}}$ of section, so $A = 40t$ and
$L = \sqrt{2A/0.175}$: \qty{21.4}{km} after \qty{1}{Myr}, \qty{42.8}{km}
after \qty{4}{Myr}, \qty{67.6}{km} after \qty{10}{Myr}. $L$ grows as
$\sqrt t$: four times as long gives twice the width. (c) $A$ is halved and $L$
divided by $\sqrt2$: \qty{15.1}{km}, \qty{30.2}{km} and \qty{47.8}{km}.
\end{solution}

\begin{solution}{exo:b1:active-margins:12}
(a) $\partial F/\partial a = -2\sum(z_i - a - s x_i) = 0$ gives
$n\bar z = na + sn\bar x$, that is $\bar z = a + s\bar x$. (b) Every $z_i$ and
$\bar z$ grow by \qty{10}{km}, so the deviations $z_i - \bar z$ and $S_{xz}$
are unchanged: the slope is the same, and $a = \bar z - s\bar x$ grows by
\qty{10}{km}. (c) The deviations $x_i - \bar x$ are unchanged, so the slope is
the same; $\bar x$ decreases by 20, so $a$ increases by $20s$. (d) Neither:
the \omterm{def:b1:geological-maps:attitude}{dip} depends only on the slope. Errors shared by all the points move the
line without turning it.
\end{solution}

\begin{solution}{pb:b1:active-margins:1}
\textbf{1.} The eight points lie close to a straight line dipping about
\ang{45}, from \qty{96}{km} deep at \qty{155}{km} from the trench to
\qty{573}{km} deep at \qty{609}{km}.

\textbf{2.} $\sum x_i = 3152$, $\sum z_i = 2751$, $\sum x_i^2 = \num{1442910}$,
$\sum x_i z_i = \num{1290548}$; $\bar x = \qty{394.0}{km}$,
$\bar z = \qty{343.875}{km}$.

\textbf{3.} $S_{xx} = 1442910 - 8\times394^2 = \num{201022}$ and
$S_{xz} = 1290548 - 8\times394\times343.875 = \num{206654}$. Then
$s = 206654/201022 = 1.028$ and $a = 343.875 - 1.02802\times394 =
\qty{-61.2}{km}$: $z = -61.2 + 1.028\,x$.

\textbf{4.} $\arctan1.028 = \ang{45.8}$.

\textbf{5.} At $x = 61.2/1.028 = \qty{59.5}{km}$ landward of the trench. A slab
that left the trench at this \omterm{def:b1:geological-maps:attitude}{dip} would reach the surface there; it reaches it at
the trench, so it \omterm{def:b1:geological-maps:attitude}{dips} more gently near the trench and steepens below.

\textbf{6.} The misfits are $-2.1$, $2.0$, $-7.3$, $5.3$, $7.1$, $2.3$,
$-15.0$ and $8.1$ \unit{km}. Event 7 lies \qty{15}{km} above the line, event 8
\qty{8}{km} below it; the others lie within about \qty{7}{km} of it, a fraction
of the thickness of the slab.

\textbf{7.} $\ang{320.9} - \ang{300} = \ang{20.9}$. Perpendicular to the
trench: $87.4\cos\ang{20.9} = \qty{81.6}{mm.yr^{-1}}$; parallel:
$87.4\sin\ang{20.9} = \qty{31.2}{mm.yr^{-1}}$. The convergence is oblique.

\textbf{8.} \qty{81.6}{mm.yr^{-1}} is \qty{81.6}{km.Myr^{-1}}: \qty{816}{km} in
\qty{10}{Myr}.

\textbf{9.} The path is $632/\sin\ang{45} = \qty{894}{km}$, covered in
$894/81.6 = \qty{11.0}{Myr}$.

\textbf{10.} $122 + 11 = \qty{133}{Myr}$.

\textbf{11.} $t = \qty{10.95}{Myr} = \qty{3.46e14}{s}$ and
$\sqrt{\num{e-6}\times\num{3.46e14}} = \qty{1.86e4}{m}$, about \qty{19}{km}
from each face: far less than the \qty{50}{km} to the middle of a slab
\qty{100}{km} thick, whose core is still cold.

\textbf{12.} $q/k$: \qty{17.2}{K.km^{-1}} beneath the old ocean floor,
\qty{14.0}{K.km^{-1}} beneath the \omterm{def:b1:active-margins:active-margin}{forearc} and \qty{32.0}{K.km^{-1}} beneath the
back-arc.

\textbf{13.} Under the \omterm{def:b1:active-margins:active-margin}{forearc} the cold slab lies only a few tens of kilometres
down; it takes up heat from the rock above it, so less heat reaches the
surface.

\textbf{14.} $h = \qty{2.60e-3}{m.s^{-2}}/(2\pi\times\num{6.674e-11}\times1700)
= \qty{3647}{m}$, about \qty{3.6}{km}: as much as, or a little more than, the
whole depth of the trench below the floor around it, about \qty{3.5}{km}. The
trench is missing mass that nothing compensates.

\textbf{15.} The density difference is $3300 - 1000 = \qty{2300}{kg.m^{-3}}$,
and $2\pi\times\num{6.674e-11}\times2300\times3500 = \qty{3.38e-3}{m.s^{-2}}$,
\qty{338}{\milli\gal}. The slab formula counts the missing mass out to any
distance; the real trench is bent down only over a band tens of kilometres
wide, so its deficit is smaller, about \qty{260}{\milli\gal}, between the two
slab estimates.

\textbf{16.} The strength of the bending plate, which holds the trench down and
the outer rise up as an elastic plate would, and the pull of the dense slab
sinking into the \omterm{def:b1:layered-earth:layers}{mantle}.

\textbf{17.} $-61.2 + 1.028\times164 = \qty{107.4}{km}$.

\textbf{18.} Yes: \qty{107}{km} lies between 72 and \qty{173}{km}, and only
about \qty{2}{km} deeper than the average, \qty{105}{km}.

\textbf{19.} $x = (100 + 61.2)/1.028 = \qty{156.8}{km}$.

\textbf{20.} $\ang{45.8}/2 = \ang{22.9}$, and
$6371\times22.9\times\pi/180 = \qty{2546}{km}$.

\textbf{21.} The fit, \ang{21.9}, is \ang{1.0} smaller, about
\qty{4}{\percent}: the bent shell describes the curvature of the Kuril arc
well.

\textbf{22.} $x = (105 + 61.2)/1.028 = 166.2/1.028 = \qty{161.7}{km}$.

\textbf{23.} $164 - 161.7 = \qty{2.3}{km}$ nearer the trench than the observed
front, $2.3/164 = \qty{1.4}{\percent}$: the Kuril volcanoes stand above a slab
very slightly deeper than the average, as question 18 found.

\textbf{24.} To the nearest kilometre, the Kuril slab line predicts an
\omterm{def:b1:active-margins:volcanic-arc}{arc--trench gap} of \textbf{\qty[text-series-to-math=true]{162}{km}} for
volcanoes above the slab at \qty{105}{km} (observed: \qty{164}{km}).
\end{solution}
